\documentclass[10pt,a4paper]{amsart}
\usepackage[margin=19mm]{geometry}
\usepackage{amsmath,amssymb,mathtools}
\usepackage[scr=rsfs,cal=euler]{mathalfa}
\usepackage{xfrac}
\usepackage[unicode,hidelinks,bookmarks=false]{hyperref}
\newcommand{\ie}{i.e.\ }
\newcommand{\cf}{cf.\ }
\newcommand{\resp}{respectively\ }
\newcommand{\Subsection}{\subsection}
\providecommand{\nobreakdash}{\nobreak\hbox{-}}
\setlength{\emergencystretch}{2em}
\multlinegap0pt
\usepackage{bm}
\usepackage{bbm}
\usepackage{dsfont}
\usepackage{xspace}
\usepackage{cancel}
\usepackage{mathtools}
\usepackage{amsthm}
\makeatletter
\@ifundefined{theorem}{\newtheorem{theorem}{Theorem}}{}
\@ifundefined{lemma}{\newtheorem{lemma}[theorem]{Lemma}}{}
\makeatother
\newcommand{\K}{\mathbb K}
\newcommand{\Z}{\mathbb Z}
\title[Path homology of circulant digraphs]{Path homology of circulant digraphs}
\author{Xinxing Tang, Shing-Tung Yau}
\date{Source: arXiv:2602.04140v1 (CC BY 4.0)}
\begin{document}
\maketitle

\noindent\emph{Source and licence.} Path homology of circulant digraphs. Version arXiv:2602.04140v1 (CC BY 4.0), distributed under the Creative Commons Attribution 4.0 International licence, as recorded in the arXiv licence metadata for this exact version and shown on the abstract page https://arxiv.org/abs/2602.04140v1. Full rights notice: licenses/arxiv-2602.04140-CC-BY-4.0.txt.\par\smallskip

\noindent\textit{Editorial scope.} Counted unit: the Theorem whose source label is \texttt{thm:S12d} in the article (source0.tex lines 800--807, source label \texttt{thm:S12d}), delivered together with the article's own complete original proof of it (source0.tex lines 809--1009, one \texttt{proof} environment of 201 lines). No mathematics is written, repaired or re-derived: statement and proof are the article's own text line for line. Required context retained uncounted: none. Every definition, display and label of those lines is retained unchanged. Omitted from this package and not redistributed: 1--799, 808--808, 1010--1948. Nothing inside a retained excerpt is omitted or altered: every retained body is delivered exactly as the article's lines are written, so no omission receipt of any kind accompanies this package. Citations: the retained excerpts carry no citation command at all, so no bibliography entry of the article is transcribed and no reference is redistributed. Reference adaptation: none was needed and none was made; every reference inside the delivered unit resolves inside it, so no unresolved reference or citation is produced. Printed slips kept literal: no printed word, symbol, hypothesis or number of the counted statement or of its proof was altered. Layout and class: the article's own class is \texttt{amsart} with packages including amssymb, amsfonts, graphicx, epstopdf, dcolumn, amsmath, enumerate, latexsym, bm, slashed, float, cite, CJK, xy; the wrapper is the lane's pinned \texttt{amsart} editorial wrapper at 10pt on A4, which restates the theorem-like environments the retained text needs and adds the article's own macro definitions that the retained text needs (\texttt{\textbackslash K}, \texttt{\textbackslash Z}), with extra wrapper packages (bm, bbm, dsfont, xspace, cancel, mathtools). The delivered numbering is the unit's own. Assets: no retained span contains a figure environment, a graphics-inclusion command, a PDF-page inclusion, a table or a TikZ picture; no image file is copied into this package, the package declares no asset dependency and no image file is redistributed with it. Licence: version arXiv:2602.04140v1 of this article is distributed under the Creative Commons Attribution 4.0 International licence, as recorded in the arXiv licence metadata for this exact version and shown on the abstract page https://arxiv.org/abs/2602.04140v1. A full rights notice is delivered at licenses/arxiv-2602.04140-CC-BY-4.0.txt. Characters: every retained excerpt is pure ASCII, so the delivered engine, XeLaTeX with its default Latin Modern text font, reports no missing character.
\begin{theorem}\label{thm:S12d} 
For $S=\{1,2,\ldots,d\}$, $2<d<n/2$. Let $i$ be the above chain map.
Then there exists a chain map
\[
\Pi:\ \left(\Omega_*(\vec{C}_n^S),\partial\right)~\rightarrow~ \left(\Omega_*(\vec{C}_n^{1,2}),\partial\right)
\]
such that $\Pi\circ i=\mathrm{Id}_{\Omega_*(\vec{C}_n^{1,2})}$ and $i\circ\Pi$ is chain-homotopic to $\mathrm{Id}_{\Omega_*(\vec{C}_n^S)}$.
\end{theorem}

\begin{proof} 
First we write an allowed elementary $m$-path in $\vec{C}_n^S$ as
\[
p=e_{v_0v_1\cdots v_m}\equiv (v_0;~ s_1,\dots,s_m),
\qquad s_r:=v_r-v_{r-1}\in\{1,\dots,d\}.
\]

\medskip
Define the \emph{long-step defect} of $p$ by
\[
E(p):=(e_1,\dots,e_m),\qquad e_r:=\max\{0,s_r-2\}\in\{0,1,\dots,d-2\}.
\]
Thus $E(p)=(0,\dots,0)$ iff every step has length $1$ or $2$, i.e.\ $p$ is an allowed $m$-path in $\vec{C}_n^{1,2}$.

We order defect profiles lexicographically:
\[
E(p)<_{\mathrm{lex}}E(q)\quad\Longleftrightarrow\quad
\exists\,r\ \text{s.t. } e_1=\cdots=e_{r-1}=e'_1=\cdots=e'_{r-1}\ \text{and } e_r<e'_r.
\]
For a finite chain $u=\sum c_p\,p\in\Omega_m(\vec{C}_n^S)$, we write
\[
E_{\max}(u):=\max\nolimits_{<_{\mathrm{lex}}}\{\,E(p): c_p\neq 0\,\}.
\]
This is a well-defined element of $\{0,\dots,d-2\}^m$.

\medskip
\noindent\textbf{Step 1: the homotopy operator $h$ and the restriction lemma.}
Define $h$ on elementary allowed $m$-paths $p=(v_0;~s_1,\dots,s_m)$ as follows.
If $E(p)=(0,\dots,0)$, set $h(p)=0$.  Otherwise, let
\[
j=j(p):=\min\{\,r\in\{1,\dots,m\}:\ s_r\ge 3\,\}.
\]
Write $s_j=\ell\ge 3$ and set
\[
h(p):=(-1)^{j+1}\,(v_0;\ s_1,\dots,s_{j-1},\,1,\,\ell-1,\,s_{j+1},\dots,s_m)\in A_{m+1}(\vec{C}_n^S).
\]
Equivalently, in vertex notation, we insert the intermediate vertex $v_{j-1}+1$ between $v_{j-1}$ and $v_j$.
Since $S=\{1,\dots,d\}$ and $\ell-1\in\{2,\dots,d-1\}$, $h(p)$ is still allowed.

Extend $h$ $\K$-linearly to $A_*(\vec{C}_n^S)$, and define
\[
\pi:=\mathrm{Id}-(\partial h+h\partial).
\]
On the full regular path space $\Lambda_*(\Z_n;\K)$ one has $\partial^2=0$, hence $\pi$ is a chain map:
$\partial\pi=\pi\partial$.

\begin{lemma}[Restriction lemma]
For every $m\ge 0$, $\pi$ preserves the GLMY subspaces:
\[
\pi\bigl(\Omega_m(\vec{C}_n^S)\bigr)\subset \Omega_m(\vec{C}_n^S).
\]
In particular, $\pi$ is a chain map on $(\Omega_*(\vec{C}_n^S),\partial)$ and $\pi|_{\Omega_*(\vec{C}_n^{1,2})}=\mathrm{Id}$.
\end{lemma}

\begin{proof}
Take $u\in\Omega_m(\vec{C}_n^S)$.  By definition, $u\in A_m(\vec{C}_n^S)$ and $\partial u\in \Omega_{m-1}(\vec{C}_n^S)\subset A_{m-1}(\vec{C}_n^S)$.
We must show $\pi(u)\in A_m(\vec{C}_n^S)$ (allowedness) and then $\partial(\pi(u))\in \Omega_{m-1}$ follows from
$\partial\pi=\pi\partial$.

\medskip\noindent
\textbf{(i) Allowedness of $h\partial u$.}
Since $\partial u\in A_{m-1}$ and $h$ sends allowed elementary paths to allowed ones (it only inserts an intermediate
vertex along a step in $S$), we have $h(\partial u)\in A_m$.

\medskip\noindent
\textbf{(ii) Allowedness of $\partial h(u)$.}
Write $u=\sum c_p p$ as a linear combination of allowed elementary $m$-paths $p$.  Then
\[
\partial h(u)=\sum c_p\,\partial h(p).
\]
We claim that every non-allowed elementary $m$-path $q$ has zero coefficient in $\partial h(u)$.

Fix such a non-allowed $q$.  Since $h(p)$ is allowed, $q$ can only appear in $\partial h(p)$ as an interior face of $h(p)$, i.e.\ by deleting an interior vertex of $h(p)$ and merging two consecutive steps whose sum exceeds $d$.
Let $j=j(p)$ be the split index for $p$, and let the inserted vertex in $h(p)$ sit between the steps
$s_{j}$ (split into $1$ and $s_j-1$).
There are three types of interior deletions in $h(p)$:

\smallskip
\emph{Type 1: deletions away from the inserted vertex.}
If we delete an interior vertex not adjacent to the inserted one, the two merged increments in $h(p)$
are identical to a pair of adjacent increments in $p$ (with the same surrounding context).
Therefore the resulting merged increment in $q$ (which is $>d$) corresponds to a forbidden $(m-1)$-face of $p$
appearing in $\partial p$.

\smallskip
\emph{Type 2: delete the vertex immediately left of the inserted vertex.}
This merges $(s_{j-1},1)$ into $s_{j-1}+1$.  If this is $>d$, then necessarily $s_{j-1}=d$, and deleting the
corresponding vertex in $p$ merges $(s_{j-1},s_j)$ into $d+s_j>d$, hence again produces a forbidden face in $\partial p$.

\smallskip
\emph{Type 3: delete the vertex immediately right of the inserted vertex.}
This merges $(s_j-1,s_{j+1})$ into $(s_j-1)+s_{j+1}$.  If this exceeds $d$, then $s_j+s_{j+1}>d+1>d$, so deleting
the corresponding vertex in $p$ merges $(s_j,s_{j+1})$ into $s_j+s_{j+1}$, again a forbidden face in $\partial p$.

\medskip
In all cases, a forbidden $m$-face $q$ of $h(p)$ canonically determines a forbidden $(m-1)$-face $F(q)$ of $p$,
obtained by also deleting the inserted vertex (i.e.\ collapsing the split) after the interior deletion.
Moreover, for fixed $q$ the set of $p$ for which $q$ occurs in $\partial h(p)$ is exactly the set of $p$ for which
$F(q)$ occurs in $\partial p$ with the same outside context; the only difference is a global sign depending on the
relative position of the inserted vertex.  Concretely, one checks in the three cases above that the coefficient
contributed by a given $p$ satisfies
\[
\mathrm{coeff}_q(\partial h(p))=\pm\,\mathrm{coeff}_{F(q)}(\partial p),
\]
where the sign $\pm$ depends only on the deletion type (hence on $q$), not on $p$.

Therefore,
\[
\mathrm{coeff}_q(\partial h(u))
=\sum_{p} c_p\,\mathrm{coeff}_q(\partial h(p))
=\pm\sum_{p} c_p\,\mathrm{coeff}_{F(q)}(\partial p)
=\pm\,\mathrm{coeff}_{F(q)}(\partial u).
\]
But $u\in\Omega_m$ implies $\partial u\in A_{m-1}$, so every forbidden $(m-1)$-path has coefficient $0$ in $\partial u$.
In particular, $\mathrm{coeff}_{F(q)}(\partial u)=0$, hence $\mathrm{coeff}_q(\partial h(u))=0$.
This holds for every forbidden $q$, so $\partial h(u)\in A_m$.

\medskip
Combining (i) and (ii), we have
\[
\pi(u)=u-\partial h(u)-h\partial u\in A_m(\vec{C}_n^S).
\]
Similarly, $\partial u\in\Omega_{m-1}$, we also have
\[\partial(\pi(u))=\pi(\partial u)\in A_{m-1}.\] 
Hence $\pi(u)\in\Omega_m$.
The identity $\pi|_{\Omega_*(\vec{C}_n^{1,2})}=\mathrm{Id}$ follows from $h=0$ on defect-zero paths.
\end{proof}

\noindent\textbf{Step 2: chain-level decrease of the defect profile.}

\begin{lemma}[Lexicographic decrease on $\Omega_m$]
Let $m\ge 0$ and $u\in\Omega_m(\vec{C}_n^S)$.  If $E_{\max}(u)\neq (0,\dots,0)$, then
\[
E_{\max}\bigl(\pi(u)\bigr) <_{\mathrm{lex}} E_{\max}(u).
\]
Consequently, for every $u\in\Omega_m$ the sequence $u,\pi(u),\pi^2(u),\dots$ stabilizes after finitely many steps
to an element with defect profile $(0,\dots,0)$, i.e.\ supported on $\{1,2\}$-steps.
\end{lemma}

\begin{proof}
Write $u=\sum c_p p$ and let $E_{\max}(u)=E(p_0)$ for some $p_0$ in the support; choose $p_0$ so that
$E(p_0)$ is maximal and among those, the first long-step index $j(p_0)$ is minimal.

Consider $\pi(p_0)=p_0-(\partial h+h\partial)(p_0)$.  By construction, the face of $h(p_0)$ obtained by deleting the inserted
vertex is exactly $p_0$, and it appears in $\partial h(p_0)$ with coefficient $+1$ (this is why we chose the prefactor
$(-1)^{j+1}$).  Hence the term $p_0$ cancels in $\pi(p_0)$.

Now examine any elementary path $q$ that appears in $\pi(p_0)$ with nonzero coefficient and is allowed
(so $q$ can contribute to $E_{\max}(\pi(u))$).  Such a $q$ arises from:
\begin{enumerate}
\item an allowed face of $h(p_0)$ other than the one deleting the inserted vertex, or
\item an allowed term in $h(\partial p_0)$ (here we use that in $\Omega_m$ all forbidden faces cancel, so only allowed faces
      survive when we pass from $\partial p_0$ to $\partial u$ and then apply $h$).
\end{enumerate}
In either case, one checks directly on step-words that $q$ is obtained from $p_0$ by one of the following local moves:
\begin{itemize}
\item either the first long step $s_{j(p_0)}=\ell$ is replaced by $\ell-1$ (hence $e_{j(p_0)}$ decreases by $1$), or
\item the first long step is moved to the right (so the defect profile agrees up to a longer prefix of zeros/unchanged
      entries, hence becomes smaller in lex order once the split has propagated).
\end{itemize}
The only way to produce a term with the same defect profile as $p_0$ would be to create a new long step strictly
to the left of $j(p_0)$ by merging shorter steps inside $\partial p_0$ and then re-splitting via $h$.
But such contributions assemble into coefficients of forbidden faces in $\partial u$ (precisely the ``partial invariance''
issue noted in blue), and they vanish because $u\in\Omega_m$ implies $\partial u$ has no forbidden faces.
Therefore, among the allowed terms that survive in $\pi(u)$, none can have defect profile $\ge_{\mathrm{lex}}E(p_0)$,
and at least one strict decrease occurs because $p_0$ itself is cancelled.  This yields the claimed inequality
$E_{\max}(\pi(u))<_{\mathrm{lex}}E_{\max}(u)$.

Since the set of defect profiles $\{0,\dots,d-2\}^m$ is finite and $<_{\mathrm{lex}}$ is a well-order on it,
iterating $\pi$ must terminate at defect $(0,\dots,0)$.
\end{proof}

\textbf{Step 3: definition of $\Pi$ and the deformation retraction.}
For $u\in\Omega_m(\vec{C}_n^S)$, define $N(u)$ to be the number of strict decreases needed to reach defect $0$, i.e.
the smallest $N$ such that $E_{\max}(\pi^N(u))=(0,\dots,0)$. Define
\[
\Pi(u):=\pi^{N(u)}(u)\in \Omega_m(\vec{C}_n^{1,2}).
\]
Because $\pi$ is a chain map on $\Omega_*(\vec{C}_n^S)$ and $\Omega_*(\vec{C}_n^{1,2})$ is a subcomplex,
$\Pi$ is a chain map $\Omega_*(\vec{C}_n^S)\to \Omega_*(\vec{C}_n^{1,2})$.
Moreover, for $u\in\Omega_*(\vec{C}_n^{1,2})$ we have $h(u)=0$ hence $\pi(u)=u$ and thus $\Pi(u)=u$; therefore
$\Pi\circ i=\mathrm{Id}$.

\medskip
\noindent\textbf{Step 4: chain homotopy between $i\circ\Pi$ and $\mathrm{Id}$.}
Define, for $u\in\Omega_m(\vec{C}_n^S)$,
\[
H(u):=\sum_{k=0}^{N(u)-1} h\bigl(\pi^k(u)\bigr)\in \Omega_{m+1}(\vec{C}_n^S),
\]
where the sum is empty (hence $H(u)=0$) if $N(u)=0$.

Using $\pi=\mathrm{Id}-(\partial h+h\partial)$ and a telescoping sum,
\[
u-\pi^{N(u)}(u)=(\partial h+h\partial)\sum_{k=0}^{N(u)-1}\pi^k(u)=\partial H(u)+H(\partial u).
\]
Since $\Pi(u)=\pi^{N(u)}(u)$ and $i$ is inclusion, this identity reads
\[
\mathrm{Id}-i\circ \Pi=\partial H+H\partial,
\]
which exhibits $i\circ\Pi$ as chain-homotopic to the identity.  This completes the proof of Theorem~\ref{thm:S12d}.
\end{proof}

\end{document}
